
\subsubsection{2.1 Training Dynamics and Per-Sample Learning}

Research on per-sample training dynamics has established that individual examples exhibit distinct learning patterns. Toneva et al. (2018) introduced the concept of forgetting events---transitions where correctly classified examples become misclassified between epochs---and showed that certain examples are forgotten repeatedly while others remain learned throughout training. Li et al. (2025) extract 142-dimensional training dynamics features per sample and demonstrate their utility for noisy label detection. Chen et al. (2025) track per-sample privacy vulnerability throughout training, finding correlations between learning difficulty and membership inference risk. Leybzon and Kervadec (2024) study memorization dynamics in language models, observing that examples memorized early are more likely to remain retained.

This line of work establishes that per-sample training dynamics carry information beyond aggregate metrics, but does not specifically address whether such dynamics can identify samples affected by spurious correlations.

\subsubsection{2.2 Spurious Correlations and Group Robustness}

Spurious correlations cause models to exploit statistical shortcuts that do not generalize \citep{Geirhos2020Shortcut}. GroupDRO \citep{Sagawa2020Distributionally} addresses this by minimizing worst-group loss, but requires group annotations. Methods reducing annotation requirements include Just Train Twice \citep{Liu2021Just}, which identifies high-error samples after initial training; Spread Spurious Attribute \citep{Nam2022Spread}, which uses pseudo-attribute prediction; and GIC \citep{Han2024Improving}, which improves group inference accuracy.

These methods focus on intervention rather than diagnosis, and rely on single-epoch signals or accuracy metrics rather than temporal loss patterns. No prior work has examined whether loss trajectory divergence between groups can specifically identify spurious correlation-affected samples.

\subsubsection{2.3 Gradient-Based Sample Analysis}

Gradient norms have been used for sample difficulty estimation \citep{Katharopoulos2018Not} and importance sampling \citep{Johnson2018Training}. Prior exploratory work found that gradient norms can identify minority samples with AUC = 0.914 on Waterbirds, though this detection did not translate into successful intervention. Maini et al. (2023) use gradient accounting to study memorization localization. These approaches provide complementary perspectives but do not specifically target spurious correlation signatures through temporal analysis.
